\section{Introduction}
\label{introduction}
The Lonely Runner conjecture was first stated by J. M. Wills in \cite{JMWills}, where it was presented as a purely formal problem in the field of Diophantine approximation. It remained so until 1998, where it was given its more memorable name, and intuitive interpretation, from \cite{Bienia97flows.view-obstructions}. As of 2010 the Lonely Runner conjecture has only been proved for up to $n$ = 6 \cite{serra_thelonely}.

There exist several equivalent interpretations of the Lonely Runner conjecture, and I will here state widely used versions\footnote{See \cite{ANote}}:
\begin{enumerate}
\item Imagine $n$ runners at a circular, unit length track\footnote{i.e. the track has a circumference of 1}, where every runner runs with constant, pairwise different speeds\footnote{These speeds can, without loss of generality be assumed to be in $\N$ - see \cite{Bienia97flows.view-obstructions} and Section \ref{integerSpeeds}}, there will then exist a time t where all runners will be at least $\frac{1}{n + 1}$ units away from their common start-point.\\

\item Imagine we have $n$ + 1 runners, then the conjecture states that for every runner there exists a time where all other runners are at least $\frac{1}{n + 1}$ units away from it. Everything else is the same as in the first interpretation \cite{Bienia97flows.view-obstructions}.\\
\end{enumerate}

More formally the problem can be stated thus: 
Given any $n$ positive integers $w_1, w_2, \ldots, w_n$, there exists a real number $x$ such that 
\eqn{
\label{eqa:lonelyRunner} \Vert w_i x\Vert \geq \frac{1}{n+1}
}

for $i = 1, 2, \ldots, n$, where for $y \in \R$, $\Vert y \Vert$ is the distance from $y$ to the nearest integer. I owe the above formulations of the conjecture to \cite{ANote}. In this report it is useful to think of the track as the one-dimensional Torus $\mathbb{T} = [0,1)$.

\subsection{Expectations of the Reader}
I expect the reader to have the mathematical knowledge required to understand both the formal and informal description of the Lonely Runner Conjecture, as well having an understanding of university level Computational Geometry\footnote{For instance having read \cite{citeulike:3347056} or similar}.

\subsection{Scope and Limitations}
\begin{itemize}
\item In this project I will not attempt to prove the Lonely Runner conjecture for all n. 
\item I will create a program that can verify whether a given configuration of runners is true for equation \eqaref{eqa:lonelyRunner}, and make it return a result, within a reasonable time, for up to 1000 runners. The final program will work with more than 1000 runners, but I will prioritise readability, maintainability and portability of the program instead of extreme optimisations. 
\item The program will include a system to test a range of runner configurations against equation \eqaref{eqa:lonelyRunner}.
\item I will only focus on runners with integer speeds - see \ref{integerSpeeds} on p. \pageref{integerSpeeds} for an in depth justification for this.
\end{itemize}

\subsection{Computer Specifications}
\label{specs}
In this section I define the computer requirements for running my program. As mentioned above, the Lonely Runner conjecture has direct applications for View-Obstruction problems and colouring of regular chromatic graphs. Researchers working with in these fields have no special reason to work with exceedingly powerful computers, or computers with special hardware.
 
Thus, any algorithm I devise and implement should be able to run on a standard computer with no special hardware\footnote{Say hardware dedicated to fast floating point calculations}. I will therefore design my verification program to run on a 2 GHz computer with the recommended amount of RAM for the computer's OS. 

I reserve the right to use third-party software libraries - as long as these can be freely distributed and installed.

\subsection{Terminology}
\label{Termonolgy}
\begin{figure}[H]
  \centering
  \includegraphics[width=0.3\textwidth, angle=90]{./images/circleZoneEPS}
  \caption{\label{circleZoneImg}An illustration of the Runner track. The greyed out part is the Zone}
\end{figure}

In this report I will refer to runners as male for simplicity.

\begin{defi}
\label{def:theZone}
In this report I will refer to the track interval [$\frac{1}{n + 1}$, $\frac{n}{n+1}$] as the Zone, where $n$ is the number of runners, and a runner who is in this interval, as ``being in the Zone''. see figure \ref{circleZoneImg} for an illustration of the Zone.
\end{defi}

\begin{defi}
\label{def:config}
Configurations of runners, will in this report refer to a set S of $n \in \N$ runners, where for all runners $r, r\prime \in S$, $r \neq r\prime$, $r_{speed}, r\prime_{speed} \in \N$, and $r_{speed} \neq r\prime_{speed}$.
\end{defi}

\hide{
\begin{defi}
\label{def:fakeRunner}
When talking about the fake runner, I will be referring to a runner that is not actually in race, but rather a runner with speed 1, who I have introduced in order to have a convenient termination criterion for the Geometrical algorithm\footnote{\pref{compGeo}}. \pref{termination} for more.
\end{defi}  
}
\subsection{Canonical Example}
Through the rest of this report I will use the following configuration
of runners to give an example:
We have 4 runners with the speeds of 3, 5, 6, 8

\subsection{Background Material Used}
\label{background}
For this project I have read the following papers dealing with the Lonely Runner conjecture: \emph{``View-Obstruction Problems''}\cite{Bienia97flows.view-obstructions}, \emph{``The lonely runner problem with seven runners''} \cite{serra_thelonely}, \emph{``Regular chromatic number and the lonely runner problem''} \cite{Barajas2007479}, \emph{``View-Obstruction Problems''} \cite{springerlink:10.1007/BF01832623}, \emph{``Tight Instances of the Lonely Runner''} \cite{Goddyn96tightinstances}, \emph{``A Note on the Lonely Runner Conjecture''} \cite{ANote} and \emph{``Invisible runners in finite fields''} \cite{invis}.

I have also used the follow support literature:
\emph{``Computational Geometry''} \cite{citeulike:3347056}, \emph{``Uniform Distribution of Sequences''} \cite{uniform} and \emph{``Kalkulus''} \cite{kalkulus}.

\subsection{Source Code and Test Data}
The source code and the data used in this report (in machine readable form) can be found in my GitHub repository at \underline{https://github.com/kjmikkel/LonelyRunner}.

\subsection{Overview}
In this section I give a quick overview of what the different sections in report will cover.
\begin{description}
\item[Section \ref{introduction}] The introduction to the report, which will introduce the Lonely Runner conjecture and explain my goals
\item[Section \ref{choiceOfMethod}] A discussion of the two different approaches that can decide whether or not equation \eqaref{eqa:lonelyRunner} holds for a given configuration
\item[Section \ref{implementation}] Implementation details and an analysis of the scalability of the implementation
\item[Section \ref{section:program}] Description of the program I have created and how to use it.
\item[Section \ref{test}] Description of the tests I have performed.
\item[Section \ref{results}] Interpretation of the results of the tests.
\item[Section \ref{conclusion}] The conclusion of the report, where I will sum up my findings and discuss future work.
\end{description}
